В ходе выполнения работы была спроектирована и реализована программная система фиксации городских событий, предназначенная для жителей города Кирова и диспетчеров муниципальных служб.

Система включает три основных компонента:
\begin{enumerate}
    \item серверную часть на платформе ASP.NET Core с базой данных PostgreSQL и расширением PostGIS, обеспечивающую хранение данных, геопространственные запросы, аутентификацию и предоставление программного интерфейса;
    \item диспетчерскую веб-панель на React с интерактивной картой обращений на базе Яндекс.Карт, дашбордом со сводной статистикой, справочниками категорий проблем и муниципальных служб;
    \item кроссплатформенное мобильное приложение на React Native с инструментарием Expo, обеспечивающее подачу обращений с фотофиксацией и привязкой к геолокации, просмотр обращений на интерактивной карте, отслеживание статусов и получение уведомлений.
\end{enumerate}

В рамках работы реализованы следующие ключевые функции:
\begin{itemize}
    \item регистрация обращений о проблемах городской инфраструктуры с фотофиксацией и автоматическим определением геолокации;
    \item автоматическое обнаружение дубликатов обращений в радиусе 50~м на основе геопространственных запросов функции ST\_DWithin расширения PostGIS;
    \item маршрутизация обращений к ответственным городским службам по категории проблемы;
    \item обработка обращений диспетчером на интерактивной карте с изменением статуса и контролем исполнения;
    \item просмотр обращений и городских событий жителем на интерактивной карте мобильного приложения;
    \item уведомление пользователя об изменении статуса его обращений.
\end{itemize}

Проведено тестирование всех компонентов системы: автоматизированные тесты серверной части на платформе xUnit (модульные и интеграционные), модульные тесты мобильного приложения на фреймворке Vitest (20 тестовых файлов), тестирование REST API через Swagger UI и функциональное тестирование веб-панели диспетчера. Все тесты завершились успешно.

Поставленная цель достигнута: разработана программная система фиксации городских событий, объединяющая мобильное приложение для жителей, диспетчерскую веб-панель и серверную часть с базой данных. Задачи работы выполнены в полном объёме.

Дальнейшее развитие программной системы может включать:
\begin{itemize}
    \item внедрение push-уведомлений на уровне операционной системы мобильного устройства;
    \item расширение возможностей работы при нестабильном сетевом соединении (кэширование данных и отложенная отправка);
    \item развитие функциональности модуля городских событий (мероприятия, перекрытия дорог, плановые работы);
    \item интеграция с внешними картографическими сервисами для автоматического обратного геокодирования (определение адреса по координатам);
    \item реализация ролевой модели с несколькими уровнями доступа для диспетчеров различных служб.
\end{itemize}
